\documentclass[a4paper,12pt]{article}
\usepackage{../../../mypackages}
\usepackage{../../../macros}

\setlength{\parindent}{0pt}

\begin{document}

\title{Chapitre 1 : Fiche d'exerices - Les nombres entiers et décimaux}
\author{N. Bancel}
\date{Septembre 2025}

\maketitle

\section{Exercice 22}

\begin{figure}[H]
  \centering
  \includegraphics[width=0.8\linewidth]{exo_22.png}
  \caption{Exercice 22}
\end{figure}

Rappel (6eme) : Pour ranger des nombres décimaux, on compare d'abord la partie entière (mètres), puis les dixièmes, puis les centièmes.

\vspace{1em}

\textbf{Question a}

    \[
      76{,}03\ \mathrm{m}
      < 76{,}59\ \mathrm{m}
      < 77{,}38\ \mathrm{m}
      < 77{,}42\ \mathrm{m}
      < 77{,}92\ \mathrm{m}
      < 78{,}80\ \mathrm{m}
      < 79{,}18\ \mathrm{m}
      < 79{,}39\ \mathrm{m}
      < 79{,}97\ \mathrm{m}
      < 84{,}12\ \mathrm{m}
    \]

  \textbf{Question b} \\
  \vspace{1em}
    Distance du \(5^{\mathrm{e}}\) lancer :

    \[
      77{,}92\ \mathrm{m}
    \]

  \textbf{Question c} \\
  \vspace{1em}
  Plus grand lancer :
  
    \[
      84{,}12\ \mathrm{m} \text{effectué par Ethan Katzberg (Canada)}
    \] 


    \begin{enumerate}
  \item \textbf{On calcule chaque nombre :}
  \[
  \begin{aligned}
  Y &= \frac{35}{10} = 3{,}5 \\
  A &= 3 + \frac{5}{10} + \frac{6}{100} = 3 + 0{,}5 + 0{,}06 = 3{,}56 \\
  P &= \frac{35}{10} + \frac{6}{1000} = 3{,}5 + 0{,}006 = 3{,}506 \\
  T &= 3 + 0{,}6 = 3{,}6 \\
  E &= 30 + \frac{506}{100} = 30 + 5{,}06 = 35{,}06 \\
  H &= \frac{3}{10} + \frac{5}{100} = 0{,}3 + 0{,}05 = 0{,}35 \\
  I &= 6 + \frac{3}{100} + \frac{2}{1000}
     = 6 + 0{,}03 + 0{,}002 = 6{,}032
  \end{aligned}
  \]
  
  \item \textbf{On range dans l'ordre croissant :}
  \[
  0{,}35 \ (H)
  < 3{,}5 \ (Y)
  < 3{,}506 \ (P)
  < 3{,}56 \ (A)
  < 3{,}6 \ (T)
  < 6{,}032 \ (I)
  < 35{,}06 \ (E)
  \]
\end{enumerate}

\end{document}
